% AI tự sinh cho Ôn thi HK Lớp 10

\begin{bt}%[0H9V2-7]
% nguon: AI tự sinh (Ôn thi HK)
	Một máy bay không người lái bay theo đường thẳng từ vị trí $A(2;3)$ đến vị trí $B(10;11)$ trong hệ trục tọa độ $Oxy$ (đơn vị: km). Tại điểm $C$ trên đường bay, người ta đặt một trạm kiểm soát sao cho khoảng cách từ $C$ đến gốc tọa độ $O$ là nhỏ nhất. Tính khoảng cách nhỏ nhất đó (làm tròn đến hàng phần trăm).
	\shortans{0,71}
	\loigiai{
		Đường thẳng $AB$ đi qua $A(2;3)$ và có vectơ chỉ phương $\overrightarrow{AB}=(8;8)$, chọn $\overrightarrow{u}=(1;1)$.\\
		Phương trình tham số: $\heva{&x=2+t\\ &y=3+t}$.\\
		Điểm $C\in AB$ nên $C(2+t;3+t)$. Khoảng cách $OC=\sqrt{(2+t)^2+(3+t)^2}$.\\
		Xét $f(t)=(2+t)^2+(3+t)^2=2t^2+10t+13$. Hàm bậc hai đạt GTNN tại $t=-\dfrac{10}{4}=-2{,}5$.\\
		Khi đó $OC_{\min}=\sqrt{2\cdot 6{,}25-25+13}=\sqrt{0{,}5}\approx 0{,}71$ km.
	}
	
\end{bt}

\begin{bt}%[0H9V2-7]
% nguon: AI tự sinh (Ôn thi HK)
	Hai bạn An và Bình cùng xuất phát từ vị trí $O(0;0)$ đi đến điểm hẹn $M(6;8)$ (đơn vị: 100 m). An đi thẳng theo đường $OM$. Bình đi theo trục $Ox$ đến điểm $A$ rồi rẽ vuông góc để đến $M$. Biết tổng quãng đường hai bạn bằng nhau, tính hoành độ của điểm $A$ (làm tròn đến hàng phần trăm).
	\shortans{0}
	\loigiai{
		Quãng đường An: $OM=\sqrt{6^2+8^2}=10$.\\
		Bình đi từ $O$ đến $A(a;0)$ rồi đến $M(6;8)$: quãng đường $=a+\sqrt{(6-a)^2+64}$.\\
		Theo đề: $a+\sqrt{(6-a)^2+64}=10\Leftrightarrow \sqrt{(6-a)^2+64}=10-a$.\\
		Điều kiện $a\le 10$. Bình phương: $(6-a)^2+64=(10-a)^2$\\
		$\Leftrightarrow 36-12a+a^2+64=100-20a+a^2\Leftrightarrow 8a=0\Leftrightarrow a=0$.\\
		Vậy $A\equiv O$, hoành độ $a=0$.
	}
	
\end{bt}

\begin{bt}%[0H9V2-5]
% nguon: AI tự sinh (Ôn thi HK)
	Trong mặt phẳng $Oxy$, một người đang ở vị trí $P(3;1)$ cần đi đến con đường thẳng $d\colon x+y-6=0$ để bắt xe buýt tại điểm $Q$, sau đó đi đến điểm làm việc $R(7;5)$. Biết rằng người đó luôn chọn $Q$ trên $d$ sao cho tổng quãng đường $PQ+QR$ nhỏ nhất. Tính tổng quãng đường nhỏ nhất đó (làm tròn đến hàng phần mười).
	\shortans{5,7}
	\loigiai{
		Xét vị trí tương đối: thay $P(3;1)$: $3+1-6=-2<0$; thay $R(7;5)$: $7+5-6=6>0$. Vậy $P,R$ khác phía $d$.\\
		Do đó $PQ+QR$ nhỏ nhất khi $P,Q,R$ thẳng hàng, tức $Q$ là giao điểm của $PR$ với $d$.\\
		$\overrightarrow{PR}=(4;4)$, chọn $\overrightarrow{u}=(1;1)$. Đường $PR$: $\heva{&x=3+t\\ &y=1+t}$.\\
		Giao với $d$: $(3+t)+(1+t)-6=0\Leftrightarrow 2t=2\Leftrightarrow t=1$. Vậy $Q(4;2)$.\\
		$PQ+QR=PR=\sqrt{4^2+4^2}=4\sqrt{2}\approx 5{,}7$.
	}
	
\end{bt}

\begin{bt}%[0H9V2-4]
% nguon: AI tự sinh (Ôn thi HK)
	Trong mặt phẳng $Oxy$, một khung thép hình chữ nhật có ba đỉnh $A(1;2)$, $B(5;-1)$, $C(m;n)$ với $AB\perp BC$. Biết $C$ nằm trên đường thẳng $y=x+3$ và hoành độ $C$ lớn hơn $5$. Tìm tung độ của $C$.
	\shortans{35}
	\loigiai{
		$\overrightarrow{AB}=(4;-3)$. Vì $AB\perp BC$ nên $\overrightarrow{AB}\cdot\overrightarrow{BC}=0$.\\
		$C(m;n)$ thuộc $y=x+3\Rightarrow n=m+3$. $\overrightarrow{BC}=(m-5;n+1)=(m-5;m+4)$.\\
		Điều kiện: $4(m-5)-3(m+4)=0\Leftrightarrow 4m-20-3m-12=0\Leftrightarrow m=32$.\\
		Khi đó $n=32+3=35$ (thỏa $m>5$). Vậy tung độ $C$ là $35$.
	}
	
\end{bt}

\begin{bt}%[0H9V2-7]
% nguon: AI tự sinh (Ôn thi HK)
	Một chiếc thuyền xuất phát từ bến $A(1;2)$ chèo thẳng đến điểm $B(9;10)$ (đơn vị: km). Do dòng nước chảy, thuyền bị lệch và thực tế di chuyển theo vectơ $\overrightarrow{v}=(1;2)$ tính theo giờ. Biết vận tốc chèo của thuyền là $5$ km/h và vận tốc dòng nước là $2$ km/h theo hướng $\overrightarrow{v}$. Tính thời gian thực tế để thuyền đến đúng điểm $B$ (làm tròn đến hàng phần trăm giờ), biết vectơ vận tốc thực là $\overrightarrow{v_{ch}}+\overrightarrow{v_n}$ với $\overrightarrow{v_{ch}}$ cùng hướng $\overrightarrow{AB}$.
	\shortans{1,63}
	\loigiai{
		$\overrightarrow{AB}=(8;8)$, $AB=8\sqrt{2}$. Đơn vị $\overrightarrow{u}_{AB}=\left(\dfrac{1}{\sqrt{2}};\dfrac{1}{\sqrt{2}}\right)$.\\
		$\overrightarrow{v_{ch}}=5\overrightarrow{u}_{AB}=\left(\dfrac{5}{\sqrt{2}};\dfrac{5}{\sqrt{2}}\right)$.\\
		Đơn vị $\overrightarrow{v}$: $|\overrightarrow{v}|=\sqrt{5}$, $\overrightarrow{u_v}=\left(\dfrac{1}{\sqrt{5}};\dfrac{2}{\sqrt{5}}\right)$, $\overrightarrow{v_n}=2\overrightarrow{u_v}=\left(\dfrac{2}{\sqrt{5}};\dfrac{4}{\sqrt{5}}\right)$.\\
		Vận tốc thực $\overrightarrow{v_{th}}=\left(\dfrac{5}{\sqrt{2}}+\dfrac{2}{\sqrt{5}};\dfrac{5}{\sqrt{2}}+\dfrac{4}{\sqrt{5}}\right)$.\\
		Thời gian $t=\dfrac{AB}{|\overrightarrow{v_{th}}|}$. Tính $|\overrightarrow{v_{th}}|^2=\left(\dfrac{5}{\sqrt2}+\dfrac{2}{\sqrt5}\right)^2+\left(\dfrac{5}{\sqrt2}+\dfrac{4}{\sqrt5}\right)^2$\\
		$=\left(12{,}5+\dfrac{20}{\sqrt{10}}+\dfrac{4}{5}\right)+\left(12{,}5+\dfrac{40}{\sqrt{10}}+\dfrac{16}{5}\right)=29{,}4+\dfrac{60}{\sqrt{10}}\approx 29{,}4+18{,}97=48{,}37$.\\
		$|\overrightarrow{v_{th}}|\approx 6{,}955$. Thời gian $t=\dfrac{8\sqrt2}{6{,}955}\approx\dfrac{11{,}314}{6{,}955}\approx 1{,}63$ giờ.
	}
	
\end{bt}

\begin{bt}%[0H9V2-6]
% nguon: AI tự sinh (Ôn thi HK)
	Trong mặt phẳng $Oxy$, một robot di chuyển trên đường thẳng $\Delta\colon x-2y+4=0$. Tại mỗi vị trí $M$ trên $\Delta$, robot đo khoảng cách đến hai trạm phát sóng đặt tại $A(1;1)$ và $B(5;3)$. Tìm giá trị nhỏ nhất của tổng $MA^2+MB^2$ (làm tròn đến hàng phần mười).
	\shortans{13,6}
	\loigiai{
		Gọi $I$ là trung điểm $AB$: $I(3;2)$. Ta có $MA^2+MB^2=2MI^2+\dfrac{AB^2}{2}$.\\
		$AB^2=(5-1)^2+(3-1)^2=16+4=20$, nên $MA^2+MB^2=2MI^2+10$.\\
		Tổng nhỏ nhất khi $MI$ nhỏ nhất, tức $M$ là hình chiếu của $I$ lên $\Delta$.\\
		$MI_{\min}=d(I,\Delta)=\dfrac{|3-2\cdot 2+4|}{\sqrt{1+4}}=\dfrac{3}{\sqrt5}$.\\
		Vậy $\min(MA^2+MB^2)=2\cdot\dfrac{9}{5}+10=\dfrac{18}{5}+10=13{,}6$.
	}
	
\end{bt}

\begin{bt}%[0H9V2-7]
% nguon: AI tự sinh (Ôn thi HK)
	Một người đi xe đạp từ $A(0;0)$ đến $B(12;5)$ (đơn vị: km) trong hệ trục $Oxy$. Trên đường đi, người đó dừng nghỉ tại điểm $M$ trên đoạn $AB$ sao cho khoảng cách từ $M$ đến trạm xăng $S(4;1)$ là nhỏ nhất. Tính khoảng cách $MS$ (làm tròn đến hàng phần trăm km).
	\shortans{0,62}
	\loigiai{
		$\overrightarrow{AB}=(12;5)$, $AB=13$. Đơn vị $\overrightarrow{u}=\left(\dfrac{12}{13};\dfrac{5}{13}\right)$.\\
		$M\in AB$: $M=A+t\overrightarrow{u}=\left(\dfrac{12t}{13};\dfrac{5t}{13}\right)$, $0\le t\le 13$.\\
		$\overrightarrow{SM}=\left(\dfrac{12t}{13}-4;\dfrac{5t}{13}-1\right)$.\\
		$SM^2=\left(\dfrac{12t}{13}-4\right)^2+\left(\dfrac{5t}{13}-1\right)^2=t^2-2t\left(\dfrac{48+5}{13}\right)+17=t^2-\dfrac{106t}{13}+17$.\\
		Đỉnh parabol: $t=\dfrac{53}{13}\approx 4{,}08\in[0;13]$.\\
		$SM^2_{\min}=17-\dfrac{53^2}{13^2}=17-\dfrac{2809}{169}=\dfrac{2873-2809}{169}=\dfrac{64}{169}$.\\
		$SM_{\min}=\dfrac{8}{13}\approx 0{,}62$ km.
	}
	
\end{bt}

\begin{bt}%[0H9V2-7]
% nguon: AI tự sinh (Ôn thi HK)
	Một khu vườn hình tam giác có ba đỉnh $A(0;0)$, $B(8;0)$, $C(2;6)$ (đơn vị: 10 m). Người ta muốn đặt một giếng nước tại điểm $G$ sao cho tổng khoảng cách từ $G$ đến ba đỉnh là nhỏ nhất (điểm Fermat). Với tam giác có góc lớn nhất nhỏ hơn $120^\circ$, điểm đó là trọng tâm. Tính khoảng cách từ $G$ đến đỉnh $A$ (làm tròn đến hàng phần mười, đơn vị 10 m).
	\shortans{3,9}
	\loigiai{
		Kiểm tra góc: $AB^2=64$, $AC^2=4+36=40$, $BC^2=36+36=72$.\\
		$\cos A=\dfrac{AB^2+AC^2-BC^2}{2\cdot AB\cdot AC}=\dfrac{64+40-72}{2\sqrt{64}\sqrt{40}}=\dfrac{32}{2\cdot 8\cdot 2\sqrt{10}}=\dfrac{1}{\sqrt{10}}>0\Rightarrow A<90^\circ$.\\
		$\cos B=\dfrac{64+72-40}{2\cdot 8\cdot 6\sqrt2}=\dfrac{96}{96\sqrt2}=\dfrac{1}{\sqrt2}\Rightarrow B=45^\circ$.\\
		$C=180^\circ-A-B>45^\circ$ nhưng kiểm tra $C<120^\circ$ (vì $A+B>60^\circ$). Vậy dùng trọng tâm.\\
		$G\left(\dfrac{0+8+2}{3};\dfrac{0+0+6}{3}\right)=\left(\dfrac{10}{3};2\right)$.\\
		$GA=\sqrt{\left(\dfrac{10}{3}\right)^2+2^2}=\sqrt{\dfrac{100}{9}+4}=\sqrt{\dfrac{136}{9}}=\dfrac{2\sqrt{34}}{3}\approx 3{,}9$.
	}
	
\end{bt}

\begin{bt}%[0H9V2-7]%[ma:3f7c1a2b]
% nguon: AI tự sinh (Ôn thi HK)
	Trong mặt phẳng tọa độ $Oxy$ (đơn vị: $100$ m), một khu dân cư có ba điểm cần kết nối hạ tầng: trạm điện $A(1;2)$, trạm cấp nước $B(7;4)$ và nhà văn hóa $C(3;6)$. Người ta muốn xây một trạm trung chuyển $M$ nằm trên đoạn thẳng $AB$ sao cho khoảng cách từ $M$ đến $C$ là nhỏ nhất. Biết $M$ chia đoạn $AB$ theo tỉ số $k=\dfrac{AM}{MB}$. Tính giá trị $k$ (làm tròn đến hàng phần trăm).
	\shortans{1,00}
	\loigiai{
		Gọi $M(x;y)$ thuộc đoạn $AB$. Vì $M$ nằm trên $AB$ nên $\overrightarrow{AM}=t\overrightarrow{AB}$ với $0\le t\le 1$.\\
		Ta có $\overrightarrow{AB}=(6;2)$, $\overrightarrow{AM}=(x-1;y-2)$.\\
		Suy ra $\heva{&x-1=6t\\ &y-2=2t}\Leftrightarrow\heva{&x=1+6t\\ &y=2+2t}$.\\
		Khoảng cách $MC^2=(x-3)^2+(y-6)^2=(6t-2)^2+(2t-4)^2$.\\
		$=36t^2-24t+4+4t^2-16t+16=40t^2-40t+20$.\\
		Hàm số $f(t)=40t^2-40t+20$ đạt giá trị nhỏ nhất tại $t=\dfrac{40}{2\cdot 40}=\dfrac{1}{2}\in[0;1]$.\\
		Khi đó $AM=\dfrac{1}{2}AB$ nên $MB=\dfrac{1}{2}AB$, do đó $k=\dfrac{AM}{MB}=1$.\\
		Vậy $k=1{,}00$.
	}
	
\end{bt}

\begin{bt}%[0H9V2-7]%[ma:9d4e5c6f]
% nguon: AI tự sinh (Ôn thi HK)
	Trong mặt phẳng tọa độ $Oxy$ (đơn vị: $10$ m), một sân bay có hai đường băng được mô hình hóa bởi hai đoạn thẳng $AB$ và $CD$ với $A(0;1)$, $B(6;5)$, $C(1;7)$, $D(7;3)$. Một xe cứu hỏa xuất phát từ vị trí $P(2;2)$ cần di chuyển đến một điểm $M$ trên đường băng $AB$ để lấy nước, sau đó đến điểm $N$ trên đường băng $CD$ để dập lửa. Biết rằng $M$ là hình chiếu vuông góc của $P$ lên đường thẳng $AB$, còn $N$ là trung điểm của đoạn $CD$. Tính độ dài quãng đường $P\to M\to N$ (làm tròn đến hàng phần mười).
	\shortans{3,8}
	\loigiai{
		\textbf{Bước 1: Tìm hình chiếu $M$ của $P$ lên $AB$.}\\
		$\overrightarrow{AB}=(6;4)$, $\overrightarrow{AP}=(2;1)$.\\
		Gọi $M=A+t\overrightarrow{AB}=(6t;1+4t)$. Điều kiện $PM\perp AB$:\\
		$\overrightarrow{PM}\cdot\overrightarrow{AB}=0\Leftrightarrow (6t-2)\cdot 6+(4t-1)\cdot 4=0$\\
		$\Leftrightarrow 36t-12+16t-4=0\Leftrightarrow 52t=16\Leftrightarrow t=\dfrac{4}{13}$.\\
		Vậy $M\left(\dfrac{24}{13};\dfrac{29}{13}\right)$.\\
		$PM=\sqrt{\left(\dfrac{24}{13}-2\right)^2+\left(\dfrac{29}{13}-2\right)^2}=\sqrt{\left(\dfrac{-2}{13}\right)^2+\left(\dfrac{3}{13}\right)^2}=\dfrac{\sqrt{13}}{13}\approx 0{,}277$.\\
		\textbf{Bước 2: Tìm trung điểm $N$ của $CD$.}\\
		$N\left(\dfrac{1+7}{2};\dfrac{7+3}{2}\right)=(4;5)$.\\
		\textbf{Bước 3: Tính $MN$.}\\
		$MN=\sqrt{\left(4-\dfrac{24}{13}\right)^2+\left(5-\dfrac{29}{13}\right)^2}=\sqrt{\left(\dfrac{28}{13}\right)^2+\left(\dfrac{36}{13}\right)^2}=\dfrac{\sqrt{784+1296}}{13}=\dfrac{\sqrt{2080}}{13}\approx 3{,}508$.\\
		\textbf{Bước 4: Tính tổng quãng đường.}\\
		$PM+MN\approx 0{,}277+3{,}508=3{,}785\approx 3{,}8$ (đơn vị $10$ m, tức $38$ m).\\
		Vậy quãng đường cần tính xấp xỉ $3{,}8$ (đơn vị tọa độ).
	}
	
\end{bt}
